%! Simplifying Algebraic Expressions and Exponent Rules — Unit 1, Lesson 1.1 — Guided Notes
\documentclass[10pt]{article}
\usepackage{atda-article}
\usepackage{atda-boxes}
\newcommand{\TallMath}[1]{$\displaystyle #1\rule[-1.4em]{0pt}{3.2em}$}
\setlength{\workrowsep}{6pt}

\begin{document}

\pageheader{Unit 1, Lesson 1.1}{Guided Notes}
% NAMESTRIP: no \namedateperiod here. The student writes their name once, on the
% cover this component is stapled behind. See references/conventions.md.

\vspace{0.15in}

\begin{objectivebox}
\small
By the end of this lesson, I will be able to\dots
\begin{enumerate}[leftmargin=1.5em, itemsep=2pt, topsep=3pt]
  \item simplify a monomial expression in one or two variables using the \textbf{exponent
        rules}: product, quotient, power of a power, power of a product, and the zero and
        negative exponents;
  \item add and subtract polynomials by combining like terms, distributing the subtraction to
        \emph{every} term of the second polynomial;
  \item multiply polynomials --- monomial by polynomial, binomial by binomial in one or two
        variables, and binomial by trinomial --- and write the result in standard form;
  \item build a polynomial expression from a situation and use it to answer a question about
        that situation, then say what the answer means.
\end{enumerate}
\end{objectivebox}

\vspace{0.10in}

\boxguard
\begin{vocabbox}
\small
Fill in each term as we build it together.
\termblanklong{Base and exponent}
\termblanklong{Like terms}
\termblanklong{Degree of a polynomial}
\termblanklong{Standard form}
\termblanklong{Distributive property}
\end{vocabbox}

\vspace{0.10in}

\boxguard[11]
\begin{hookbox}
\small
\textbf{The Party Size tub.} A theater's regular popcorn tub is a cylinder with radius $3$ inches
and height $8$ inches. The \emph{Party Size} is exactly twice as wide and twice as tall ---
radius $6$ inches, height $16$ inches --- and it costs \textbf{three times} as much. The volume
of a cylinder is $V=\pi r^2 h$.

Is the Party Size a good deal? Decide first, then check it.

\writelines{2}
\end{hookbox}

\vspace{0.10in}

\boxguard[26]
\begin{notesbox}{1. The Exponent Rules --- Count the Factors}
\small
An exponent counts factors. Every rule below is what happens when you write the factors out and
count them, so none of them has to be memorized cold.

\renewcommand{\arraystretch}{1.9}
\begin{tabularx}{\linewidth}{@{}p{3.6cm} p{4.0cm} X@{}}
\rowcolor{mist}
\textbf{Rule} & \textbf{In symbols} & \textbf{Try it} \\
\hline
Product of powers   & $a^m \cdot a^n = a^{m+n}$ & $x^4 \cdot x^9 =$ \blank{2.4cm} \\
Quotient of powers  & $\dfrac{a^m}{a^n} = a^{m-n}$ & $\dfrac{y^{11}}{y^4} =$ \blank{2.4cm} \\
Power of a power    & $\left(a^m\right)^n = a^{mn}$ & $\left(c^3\right)^5 =$ \blank{2.4cm} \\
Power of a product  & $(ab)^n = a^n b^n$ & $(2d)^4 =$ \blank{2.4cm} \\
Power of a quotient & $\left(\dfrac{a}{b}\right)^n = \dfrac{a^n}{b^n}$
                    & $\left(\dfrac{k}{3}\right)^2 =$ \blank{2.4cm} \\
\end{tabularx}

\vspace{6pt}
\textbf{Two worked examples.} Coefficients \emph{multiply}; exponents \emph{add}. They are
different operations and mixing them up is the most common error on this page.
\begin{work}
  \left(-2p^3q^4\right)\left(5p^2q\right) &= (-2)(5)\,p^{3+2}\,q^{4+1} \\
                                          &= -10p^5q^5
\end{work}
\begin{work}
  \left(2m^3n\right)^4 &= 2^4\left(m^3\right)^4 n^4 \\
                       &= 16m^{12}n^4
\end{work}

\textbf{Careful:} in $\left(2m^3n\right)^4$ the exponent $4$ reaches the \blank{2.8cm} as well as
the variables --- that is why the answer starts with $16$, not $2$.
\end{notesbox}

\vspace{0.10in}

\boxguard[15]
\begin{notesbox}{2. Zero and Negative Exponents --- Where They Come From}
\small
Nobody decided these; they are forced on us by the quotient rule.

\begin{itemize}[leftmargin=1.3em, itemsep=4pt, topsep=3pt]
  \item $\dfrac{x^5}{x^5}$ is $x^{5-5}=x^0$ by the quotient rule, but every factor cancels, so it
        is also $1$. Therefore $x^0=1$ for any $x \ne 0$.
  \item $\dfrac{x^3}{x^7}$ is $x^{3-7}=x^{-4}$ by the quotient rule, but cancelling three factors
        leaves $\dfrac{1}{x^4}$. Therefore $x^{-4}=\dfrac{1}{x^4}$.
\end{itemize}

\textbf{Say it correctly:} a negative exponent does not make the value \blank{3.0cm}. It means
the factor belongs on the other side of the fraction bar.

\begin{work}
  \dfrac{24a^7b^2}{6a^3b^5} &= 4\,a^{7-3}\,b^{2-5} \\
                            &= 4a^4b^{-3} \\
                            &= \dfrac{4a^4}{b^3}
\end{work}
\end{notesbox}

\vspace{0.10in}

\boxguard[16]
\begin{notesbox}{3. Adding and Subtracting Polynomials}
\small
Adding is just collecting like terms. Subtracting is the same --- \emph{after} the minus sign
reaches every term of the second polynomial.
\begin{work}
  \left(4x^2-7x+2\right)+\left(x^2+3x-9\right) &= \left(4x^2+x^2\right)+(-7x+3x)+(2-9) \\
                                               &= 5x^2-4x-7
\end{work}
\begin{work}
  \left(4x^2-7x+2\right)-\left(x^2+3x-9\right) &= 4x^2-7x+2-x^2-3x+9 \\
                                               &= 3x^2-10x+11
\end{work}

\textbf{The trap:} the subtraction sign belongs to \blank{2.8cm} term of the second polynomial,
not just the first one. Watch the $-9$ become $+9$ above.

\vspace{4pt}
\textbf{In context --- the concession stand.} At a price of $x$ dollars per nacho tray, the stand
expects to sell $150-10x$ trays. It pays \$40 for the night plus \$2 in ingredients per tray sold.

\begin{enumerate}[leftmargin=1.6em, itemsep=5pt, topsep=4pt]
  \item Revenue is price times trays sold; cost is the flat fee plus \$2 per tray. Write each in
        standard form.
\begin{work}
  R(x) &= x(150-10x) = -10x^2+150x \\
  C(x) &= 40+2(150-10x) = -20x+340
\end{work}

  \item Profit is revenue minus cost. Find $P(x)$.
\begin{work}
  P(x) &= \left(-10x^2+150x\right)-\left(-20x+340\right) \\
       &= -10x^2+150x+20x-340 \\
       &= -10x^2+170x-340
\end{work}

  \item Find $P(8)$ and write one sentence saying what it means for the stand.
\begin{work}
  P(8) &= -10(8)^2+170(8)-340 \\
       &= -640+1360-340 \\
       &= 380
\end{work}
\writeline
\end{enumerate}
\end{notesbox}

\vspace{0.10in}

\boxguard[16]
\begin{notesbox}{4. Multiplying Polynomials --- Every Term, Every Time}
\small
There is only one rule: \textbf{every term of the first factor multiplies every term of the
second}, and then you combine like terms. A $2$-term factor times a $3$-term factor produces
\blank{1.8cm} products before combining --- count them as a check.

\begin{work}
  3x^2\left(2x^2-5x+4\right) &= 3x^2\left(2x^2\right)-3x^2(5x)+3x^2(4) \\
                             &= 6x^4-15x^3+12x^2
\end{work}

\textbf{Two variables} (\texttt{A2.EO.3a} asks for these too --- $ab$ terms are like terms only
with other $ab$ terms):
\begin{work}
  (3a+2b)(a-5b) &= 3a^2-15ab+2ab-10b^2 \\
                &= 3a^2-13ab-10b^2
\end{work}

\textbf{Binomial times trinomial} --- six products, then combine:
\begin{work}
  (x+3)\left(x^2-2x+5\right) &= x^3-2x^2+5x+3x^2-6x+15 \\
                             &= x^3+x^2-x+15
\end{work}
\end{notesbox}

\vspace{0.10in}

\boxguard[16]
\begin{practicebox}
\small
A concession tray is folded from a $10$-inch by $8$-inch sheet of cardboard: cut an $x$-inch
square out of each corner, then fold the four flaps up. The volume in cubic inches is
\[ V(x) = x(10-2x)(8-2x). \]

\begin{enumerate}[leftmargin=1.6em, itemsep=6pt, topsep=4pt]
  \item Expand $V(x)$ and write it in standard form. (Multiply the two binomials first.)
\begin{work}
  (10-2x)(8-2x) &= 80-20x-16x+4x^2 \\
                &= 4x^2-36x+80 \\
  V(x) &= x\left(4x^2-36x+80\right) \\
       &= 4x^3-36x^2+80x
\end{work}

  \item Find $V(1)$, then say what that number means for the tray.
\begin{work}
  V(1) &= 4(1)^3-36(1)^2+80(1) \\
       &= 4-36+80 \\
       &= 48
\end{work}
\writeline

  \item The cut size $x$ cannot be $4$ inches or more. Use the factored form to explain why.

\writelines{2}
\end{enumerate}
\end{practicebox}

\end{document}
